\subsection{Mathematical Language, Sets, and Functions}

% **Role:** rigorous language of real and complex analysis.
%
% Logic, quantifiers, induction, proof methods, set operations, functions, relations.


This subsection fixes the basic language of Analysis: membership, containment, functions, images, preimages, and definition-level proof patterns.

\begin{principle} \label{princ:analysis-language-as-structure} \hlt{Language Carries Structure}\\
Objects, hypotheses, conclusions, and quantifier order matter. For instance,
\begin{align}
\forall x\in X\ \exists y\in Y\ P(x,y), \qquad
\exists y\in Y\ \forall x\in X\ P(x,y) \nonumber
\end{align}
are different: in the first $y$ may depend on $x$; in the second one $y$ works for all $x$.
\end{principle}

\begin{principle} \label{princ:well-formed-statements-and-bound-variables} \hlt{Well-Formed Statements and Bound Variables}\\
A mathematical statement is well formed only after every symbol has a declared role and every expression has a defined meaning. A variable introduced by a quantifier, summation sign, integral, set-builder condition, or similar construction is \hlt{bound} inside that construction; changing its name does not change the statement as long as its scope is unchanged and no free variable is captured.
\end{principle}

\begin{definition} \label{def:finite-sums-indexed-by-integers} \hlt{Finite Sums Indexed by Integers}\\
Let $a_i\in\C$ be assigned for integers $m\leq i\leq n$. Define $\sum\limits_{i=m}^{n}a_i\defeq a_m+a_{m+1}+\cdots+a_n$.\\
If $m>n$, the sum over the empty index set is $0$.
\end{definition}

\begin{proposition} \label{prop:finite-sum-basic-laws} \hlt{Basic Laws for Finite Sums}\\
Let $a_i,b_i\in\C$ and $c\in\C$ be assigned for $m\leq i\leq n$.
\begin{enumerate}[label=(\roman*)]
\setlength{\itemsep}{0pt}
\item If $m\leq r<n$, then $\sum\limits_{i=m}^{n}a_i=\sum\limits_{i=m}^{r}a_i+\sum\limits_{i=r+1}^{n}a_i$.
\item $\sum\limits_{i=m}^{n}(a_i+b_i)=\sum\limits_{i=m}^{n}a_i+\sum\limits_{i=m}^{n}b_i$ and $\sum\limits_{i=m}^{n}ca_i=c\sum\limits_{i=m}^{n}a_i$.
\item $\abs{\sum\limits_{i=m}^{n}a_i}\leq\sum\limits_{i=m}^{n}\abs{a_i}$.
\item If $a_i,b_i\in\R$ and $a_i\leq b_i$ for every $i$, then $\sum\limits_{i=m}^{n}a_i\leq\sum\limits_{i=m}^{n}b_i$.
\end{enumerate}
\begin{proof}
\begin{enumerate}[label=(\roman*)]
\setlength{\itemsep}{0pt}
\item By Definition \ref{def:finite-sums-indexed-by-integers}, associativity of addition in $\C$, $\sum\limits_{i=m}^{n}a_i=\sum\limits_{i=m}^{r}a_i+\sum\limits_{i=r+1}^{n}a_i$.
\item By associativity, commutativity, and distributivity in the field $\C$ (Theorem \ref{thm:complex-numbers-form-field}),\\
$\sum\limits_{i=m}^{n}(a_i+b_i)=\sum\limits_{i=m}^{n}a_i+\sum\limits_{i=m}^{n}b_i$ and $\sum\limits_{i=m}^{n}ca_i=c\sum\limits_{i=m}^{n}a_i$.
\item Induct on $n-m$. The case $n=m$ is immediate.\\
If the claim holds up to $n$, then by Theorem \ref{thm:complex-absolute-value-identities}, $\abs{\sum\limits_{i=m}^{n+1}a_i} \leq \abs{\sum\limits_{i=m}^{n}a_i}+\abs{a_{n+1}} \leq \sum\limits_{i=m}^{n}\abs{a_i}+\abs{a_{n+1}} =\sum\limits_{i=m}^{n+1}\abs{a_i}$.
\item Induct on $n-m$. The case $n=m$ is the hypothesis $a_m\leq b_m$.\\
If the claim holds up to $n$, then Proposition \ref{prop:ordered-field-inequality-rules-analysis} gives $\sum\limits_{i=m}^{n+1}a_i =\sum\limits_{i=m}^{n}a_i+a_{n+1} \leq \sum\limits_{i=m}^{n}b_i+b_{n+1} =\sum\limits_{i=m}^{n+1}b_i$.
\end{enumerate}
\hlt{Structural Commentary}: Finite sums are iterated addition, so their laws are inherited from the scalar field.\\
\hlt{Key Insight}: No convergence hypothesis is needed for a finite sum.
\end{proof}
\end{proposition}

\begin{principle} \label{princ:mathematical-induction} \hlt{Mathematical Induction}\\
Let $P_1,P_2,\ldots$ be statements indexed by positive integers. To prove $P_n$ for every $n\in\N$, it is enough to prove:
\begin{enumerate}[label=(\roman*)]
\setlength{\itemsep}{0pt}
\item \hlt{basis}: $P_1$ is true;
\item \hlt{induction step}: for every $n\in\N$, $P_n$ implies $P_{n+1}$.
\end{enumerate}
The induction step is an implication; it does not prove $P_{n+1}$ in isolation.\\
The basis starts the chain, and the implication carries truth from one statement to the next.
\end{principle}

\begin{figure}[H]
  \centering
  \begin{tikzpicture}[
  >=Stealth,
  every node/.style={font=\small},
  value/.style={rounded corners=2pt, draw=orange!70!black, fill=orange!10, minimum width=10mm, minimum height=7mm, inner xsep=4pt, line width=0.55pt},
  arrow/.style={->, line width=0.5pt, gray!62},
  rowlabel/.style={anchor=east, gray!68, font=\small},
  rowsub/.style={anchor=east, gray!58, font=\scriptsize}
]
  \path[use as bounding box] (-2.65,-1.85) rectangle (6.25,0.52);

  \node[rowlabel] at (-0.72,0.12) {induction};
  \node[rowsub] at (-0.72,-0.18) {base + step};
  \node[value] (p0) at (0,0) {$P(0)$};
  \node[value] (p1) at (1.55,0) {$P(1)$};
  \node[value] (p2) at (3.10,0) {$P(2)$};
  \node[value] (p3) at (4.65,0) {$P(3)$};
  \node at (5.85,0) {$\cdots$};

  \draw[arrow] (p0) -- (p1);
  \draw[arrow] (p1) -- (p2);
  \draw[arrow] (p2) -- (p3);
  \draw[arrow] (p3) -- (5.55,0);

  \node[rowlabel] at (-0.72,-1.28) {recursion};
  \node[rowsub] at (-0.72,-1.58) {initial + rule};
  \node[value] (a0) at (0,-1.40) {$a_0=c$};
  \node[value] (a1) at (1.55,-1.40) {$a_1$};
  \node[value] (a2) at (3.10,-1.40) {$a_2$};
  \node[value] (a3) at (4.65,-1.40) {$a_3$};
  \node at (5.85,-1.40) {$\cdots$};

  \draw[arrow] (a0) -- (a1);
  \draw[arrow] (a1) -- (a2);
  \draw[arrow] (a2) -- (a3);
  \draw[arrow] (a3) -- (5.55,-1.40);
\end{tikzpicture}

  \caption{Induction propagates truth, while recursion propagates values along the same index chain.}
  \label{fig:peano-recursion-induction-ladder}
\end{figure}

\begin{method} \label{meth:proof-grammar-and-negation} \hlt{Proof Grammar and Negation}\\
In theorem statements, put the hypotheses before the conclusion and keep quantifiers in the order meant.\\
In particular, switching $\forall$ and $\exists$ usually changes the statement; see Principle \ref{princ:analysis-language-as-structure}.\\
Before applying a rule, also check that the displayed expressions are well formed; an undefined expression makes the sentence ill formed, not merely false.\\
The word ``or'' is inclusive in mathematics, and the basic negation rules are
\begin{align}
\neg(a\land b)\Longleftrightarrow \neg a\lor\neg b, \qquad \neg(a\lor b)\Longleftrightarrow \neg a\land\neg b, \qquad \neg(a\Rightarrow b)&\Longleftrightarrow a\land\neg b \nonumber
\end{align}
Thus a proof by contradiction of $a\Rightarrow b$ assumes both $a$ and $\neg b$ and derives an impossibility.\\
\hlt{Key Insight}: Before proving or negating a statement, parse its quantifiers and logical connectives.
\end{method}
